\section{Conclusion}
\label{sec:conclusion}

\cyberlite{} asks whether a sub-billion-parameter language model can be
post-trained into a useful cybersecurity assistant under consumer hardware
constraints. The project contributes an architecture-aware pipeline for
Qwen3.5-0.8B, a provenance-first dataset methodology, a non-executing local
inference and autocomplete stack, and a task-specific evaluation harness with
explicit uncertainty and memory reporting.

At this stage, the honest conclusion is about readiness, not performance. The
pipeline runs end to end on CPU at smoke scale, the hybrid checkpoint loads and
trains through LoRA, adapters reload for inference, and the test suite validates
the data, metrics, retrieval, autocomplete, and API paths without downloading
model weights. The research conditions, however, have not been run, and this
manuscript reports no benchmark numbers.

The central philosophy is that cybersecurity AI should be accessible without
enterprise-grade hardware. Accessibility without rigor would be a liability.
The project therefore pairs low-resource engineering with licensed data,
contamination screening, fixed seeds, confidence intervals, measured memory, and
a refusal to invent results. If the hypotheses hold, the outcome is a practical
local assistant and a reusable research stack. If they fail, the negative result
will still clarify the capability floor of compact models under realistic
constraints.

We invite independent reproduction, criticism, and collaboration---especially
from researchers and practitioners who work outside well-resourced labs.
